%% 致谢
\chapter*{致\quad\quad 谢}\addcontentsline{toc}{chapter}{致谢}

我的导师段世红老师在选题方向的敲定、CityLearn 仿真框架的取舍、CHESCA 基线与残差多智能体改进方案的设计，以及可视化管理平台的需求分析等环节给出了系统性指导；其在第 17 天电网突发停电压力测试、跨季节边界实验这两项对工程落地说服力较强的实验设计上也给出了直接建议，使本研究能够从规则协同、可复现性与平台可运维三条线索并行推进。

北京科技大学数理学院的各位老师在课程学习、开题评审与中期检查阶段指出了关键短板，帮助我把研究边界收窄在"规则协同优先 + 纯 MARL 上限探明 + 全栈可复现平台"三件事上。同门与实验室同学在算力平台调试、异步任务调度链路 debug 与论文图表统一性校核中分担了大量工作。

最后，家人撑起了三年学习期间的生活底色，评审专家与答辩委员的意见将进一步校准本研究的应用口径。
